% coverletter.tex
% Pairs with resume.tex — same theme, same header block. Both pull
% from ../shared/theme.sty and ../shared/title.tex, so the cover
% letter and resume always stay visually aligned automatically; you
% never need to hand-sync fonts/colors/contact info between the two.
% Content is auto-generated from data/coverletter.yaml into
% sections/coverletter_content.tex by scripts/build.py — edit the
% YAML, not this file.
\documentclass[11pt]{article}

\usepackage[margin=1in]{geometry}
\usepackage{hyperref}
\usepackage{../shared/theme}      % colors/fonts, generated from data/theme.yaml
% AUTO-GENERATED by scripts/build.py from data/title.yaml — do not hand-edit.
\newcommand{\myName}{James Smith}
\newcommand{\myTagline}{Mechatronics Engineering Student}
\newcommand{\myLocation}{Waterloo, ON}
\newcommand{\myEmail}{JamesThomasSmith07@icloud.com}
\newcommand{\myPhone}{}
\newcommand{\myGithub}{https://github.com/james-t-smith}
\newcommand{\myLinkedin}{https://www.linkedin.com/in/james-smith-984b2b3a8/}
\newcommand{\myWebsite}{https://james-t-smith.github.io/}
       % \myName, \myEmail, etc. — generated from data/title.yaml

\hypersetup{colorlinks=true, urlcolor=blueprint500, linkcolor=blueprint500}
\pagestyle{empty}
\setlength{\parindent}{0pt}
\setlength{\parskip}{0pt}

\begin{document}

% ---- header — identical treatment to resume.tex's header, so the
% two documents read as one set when stapled or attached together ----
\begin{center}
  {\displayfont\bfseries\LARGE \myName}\\[2pt]
  {\color{steel700} \myTagline}\\[4pt]
  \small \myLocation\ \textbar\ \myEmail\ \textbar\ \myPhone\ \textbar\
  \href{\myGithub}{GitHub} \textbar\ \href{\myLinkedin}{LinkedIn}
\end{center}

\vspace{1.2em}
{\color{blueprint500}\hrule height 0.8pt}
\vspace{1.4em}

% ---- body — auto-generated from data/coverletter.yaml, do not hand-edit sections/coverletter_content.tex ----
% AUTO-GENERATED by scripts/build.py from data/coverletter.yaml
% Do not hand-edit — edit the YAML instead and re-run the build.

September 13, 2026

\vspace{12pt}
Dear Hiring Team,

\vspace{8pt}
I'm writing to apply for the Mechatronics Engineering Intern position at Example Robotics Co. As a mechatronics engineering student with hands-on experience across embedded systems, PCB design, and control theory, I'm excited about the opportunity to contribute to your automation team.

\vspace{8pt}
In my most recent project, I designed and tuned a PID-controlled line-following robot, improving lap times by 38\% over a fixed-gain baseline — the kind of iterative, measurement-driven engineering I'd bring to your team. I've also led a two-person hardware team building a stackable PCB sensor system, and derived closed-form inverse kinematics for a 3-DOF robotic arm.

\vspace{8pt}
I'd welcome the chance to talk about how my background fits your team's needs. Thank you for your time and consideration.

\vspace{8pt}
Sincerely,\\[28pt]
\myName


\end{document}
